% !TeX program = lualatex
\documentclass[a4paper,11pt]{ltjsarticle}
\usepackage[margin=25mm]{geometry}
\usepackage{fvextra}
\usepackage[hidelinks]{hyperref}

\setlength{\parindent}{1em}
\setlength{\parskip}{0.6em}
\fvset{
  fontsize=\tiny,
  frame=single,
  numbers=left,
  numbersep=4pt,
  xleftmargin=0.5em,
  breaklines=true,
  breakanywhere=true,
  breakautoindent=false,
  breakindent=0pt,
  breaksymbolleft={}
}

\title{生成AIとRAGを用いた観光プラン生成アプリの実装と考察\\\large Tavily・Gemini を用いた Ehime Tour Planner の開発}
\author{}
\date{2026年4月20日}

\begin{document}
\maketitle

\section*{(1) やったこと・動かしたものの概要}
本課題では，生成 AI と検索拡張生成（RAG）を組み合わせた観光プラン生成システム
\texttt{Ehime Tour Planner} を実際に実装し，ローカル環境および公開済みアプリ上で動作確認を行った。
単に既存アプリを紹介するのではなく，
\texttt{Tavily} による観光情報検索，
\texttt{gemini-embedding-001} による埋め込み表現，
\texttt{FAISS} を用いた関連文書検索，
\texttt{gemini-2.5-flash-lite} による構造化された旅程生成という一連の AI 技術を組み合わせ，
入力から出力までの流れを自分で動かして確かめた。

具体的には，利用者が旅行日数，同行者，移動手段，関心テーマ，訪問エリアなどを入力すると，
まず観光情報ページを検索・取得し，
その内容を埋め込みベクトルに変換して関連箇所を抽出し，
その後に LLM が出典 URL 付きの旅行プランを生成する仕組みを作成した。
さらに，初回プランの出力後にチャットで修正指示を与え，
既存プランを再構成できることも確認した。

この課題を通じて実際に動かした対象は，
単体の API 呼び出しではなく，
検索，埋め込み，検索結果の選別，生成，対話的な再生成を統合した小規模な生成 AI アプリケーションである。
そのため，本レポートではアプリの見た目そのものよりも，
どの AI 技術をどの順序で用い，
どのような処理として実装したかに重点を置いて述べる。

\noindent\textbf{使用技術の要点}
\begin{itemize}
\item UI: \texttt{Streamlit}
\item 検索: \texttt{Tavily Search}（必要に応じて Web 全体も追加）
\item 埋め込み: \texttt{gemini-embedding-001}
\item 生成: \texttt{gemini-2.5-flash-lite}
\item 検索基盤: \texttt{FAISS}，利用できない場合は NumPy のコサイン類似度で代替
\item 出力: JSON スキーマ準拠の旅程と Markdown ダウンロード
\end{itemize}

\section*{(2) 具体的な手続き}
\subsection*{1. 環境設定と実行手順}
まず，生成 AI アプリを実際に動作させるための環境を整えた。
本システムは Python 環境上で \texttt{Streamlit} を用いて実行した。
必要なライブラリは \texttt{requirements.txt} にまとめられており，
\texttt{streamlit}，\texttt{google-genai}，\texttt{tavily-python}，
\texttt{faiss-cpu}，\texttt{beautifulsoup4} などを利用した。

実行準備としては，まず依存関係をインストールし，
つぎに \texttt{.streamlit/secrets.toml} に
\texttt{GEMINI\_API\_KEY} と \texttt{TAVILY\_API\_KEY} を設定する。
\texttt{app.py} では，これらの値を \texttt{st.secrets} または環境変数から読み込み，
どちらかが欠けている場合はアプリを停止する実装になっている。

実行手順の例を以下に示す。
\begin{itemize}
\item \texttt{pip install -r requirements.txt}
\item \texttt{.streamlit/secrets.toml} を作成し，API キーを設定
\item \texttt{streamlit run app.py}
\end{itemize}

このようにして環境構築後，
ローカル実行に加えて公開済みの Streamlit アプリでも同じ処理が動くことを確認した。公開先は以下の通りである。\par
\url{https://ehime-tour-planner-iicqpkkbfs9zrcjba6a9at.streamlit.app/}

\subsection*{2. 入力条件を受け取る画面}
実装の最初の段階では，利用者が旅行条件を指定するための画面を \texttt{Streamlit} で構成した。
旅行日数，開始日，同行者，移動手段，発着地，関心テーマ，訪問エリア，
子連れ配慮，1 日の詰め込み度をサイドバーから入力できるようにしている。
この段階で条件を明示的に取得しておくことで，
後段の検索クエリと旅程生成プロンプトに同じ制約を反映できる。

\noindent\textbf{コード 1: 条件入力 UI（\texttt{app.py} より）}
\VerbatimInput[firstline=46,lastline=91,firstnumber=46]{../app.py}

\subsection*{3. 関連ページの収集と前処理}
次に，RAG の入力となる観光情報を集める処理を実装した。
このアプリでは \texttt{Tavily} を用い，
既定では \texttt{iyokannet.jp} にドメインを限定して検索する。
必要に応じて Web 全体の検索結果も加えられるようにしつつ，
URL の重複は除去し，
取得した本文は \texttt{BeautifulSoup} でテキスト化してから保持する。
これにより，出典を保ちつつも，
そのままの転載ではない再利用可能な検索用データを作成できる。

\noindent\textbf{コード 2: 検索と前処理（\texttt{rag/retriever.py} より）}
\VerbatimInput[firstline=38,lastline=120,firstnumber=38]{../rag/retriever.py}

\subsection*{4. 埋め込みとベクトル検索による RAG}
収集したテキストはそのまま生成器に渡さず，
一度チャンクに分割して埋め込みベクトルへ変換する。
本アプリでは \texttt{gemini-embedding-001} を使い，
\texttt{FAISS} が利用可能であればそのインデックスを用い，
利用できない環境では NumPy によるコサイン類似度計算へ切り替える。
この実装により，
環境差に強い RAG パイプラインを保ちながら，
ユーザーの検索語に近い文脈だけを抽出できる。

\noindent\textbf{コード 3: 埋め込みと検索基盤（\texttt{rag/retriever.py} より）}
\VerbatimInput[firstline=122,lastline=183,firstnumber=122]{../rag/retriever.py}

さらに \texttt{retrieve\_for\_plan} では，
単に類似度の高いチャンクを並べるのではなく，
「1 URL あたり 1 チャンク」に絞ることで情報源の偏りを抑えている。
また，選ばれたチャンク群は Gemini に一括で要約させる。
このまとめ要約は API 呼び出し回数の削減につながるだけでなく，
観光記事の原文を長く引用しないという方針にも合致している。

\noindent\textbf{コード 4: チャンク選択と一括要約（\texttt{rag/retriever.py} より）}
\VerbatimInput[firstline=194,lastline=307,firstnumber=194]{../rag/retriever.py}

\subsection*{5. プロンプト設計と構造化出力}
検索で得られた文脈をもとに旅程を生成する段階では，
自由文だけを返させるのではなく，
JSON スキーマに従った構造化出力を要求している。
プロンプト側では，
長い原文引用の禁止，
各日ごとの根拠 URL の提示，
安全面や季節性への配慮などをガードレールとして明示する。
また，発着地や宿泊地の扱いを具体的に指定することで，
単なる観光地の列挙ではなく，
日程として実行しやすいプランになるよう工夫している。

\noindent\textbf{コード 5: 旅程生成プロンプト（\texttt{rag/prompts.py} より）}
\VerbatimInput[firstline=67,lastline=127,firstnumber=67]{../rag/prompts.py}

実際の生成処理では，
上記プロンプトと抽出済みコンテキストを \texttt{generate\_content} に渡し，
\texttt{response\_json\_schema} を指定して旅程の構造を固定する。
これにより，生成結果をそのまま画面表示や Markdown 出力へつなげやすくなる。

\noindent\textbf{コード 6: 初回プラン生成（\texttt{app.py} より）}
\VerbatimInput[firstline=131,lastline=159,firstnumber=131]{../app.py}

\subsection*{6. チャットによる再計画}
本アプリでは，初回生成のあとにチャットで修正依頼を受け取り，
既存プランを保ちながら再構成する流れを用意している。
利用者は「2 日目をもっとゆったりにしたい」
「温泉を追加したい」といった自然文で修正を依頼できる。

\noindent\textbf{コード 7: チャットによるプラン修正（\texttt{app.py} より）}
\VerbatimInput[firstline=188,lastline=220,firstnumber=188]{../app.py}

\section*{(3) わかったこと・わからなかったこと}
\subsection*{1. わかったこと}
実際にシステムを動かしてみて，次の点が重要だと分かった。

\begin{itemize}
\item 検索対象を \texttt{iyokannet.jp} 中心に絞ることで，観光プランの根拠を明確にしやすい。
\item 生成前に関連チャンクを選別し，要点だけをまとめて渡すことで，長い観光記事をそのまま扱うよりも安定した旅程が得られる。
\item JSON スキーマによる構造化出力は，表示，保存，再編集のすべてに有効であり，単なるテキスト生成より実用性が高い。
\item チャットで再計画できる構成にしておくと，初回出力をたたき台として徐々に利用者の意図へ寄せられる。
\end{itemize}

\subsection*{2. わからなかったこと・この講義で解消したいこと}
一方で，実装してみてまだ十分に整理できていない点も残った。
\begin{itemize}
\item 埋め込みベクトルを用いると，なぜ意味的に近い観光文書どうしが近い位置に配置され，検索に有効になるのかをより深く理解したい。表現学習や自己教師あり学習の観点から仕組みを整理したい。
\item 今回は Gemini を用いて旅行プランを生成したが，Transformer や大規模言語モデルが，従来の系列モデルよりどのように文脈を扱い，長い入力から要点を取り出しているのかを講義で学びたい。
\item RAG を使っても，生成結果が不正確になったり，根拠の弱い内容を含んだりすることがある。深層生成モデルや自然言語処理の観点から，そのような誤りがなぜ起こるのか，またどうすれば抑えられるのかを理解したい。
\end{itemize}

\section*{参考リンク}
\begin{itemize}
\item Ehime Tour Planner アプリ公開リンク:\par
\url{https://ehime-tour-planner-iicqpkkbfs9zrcjba6a9at.streamlit.app/}
\end{itemize}

\end{document}
